\chapter{Case Study on Migrating Data Contracts to Ontologies}
\label{ch:case}

\section{Introduction}
This chapter reports a case study conducted according to the guidelines of
\citeonline{runeson2012case}. A case study fits the problem because the migration is a
contemporary phenomenon that can only be observed inside a real system: the contracts, the
server that consumes them and the questions asked by its users are not under the control of
the researcher. The study is a single-case, descriptive and exploratory study. Its purpose
is to characterise the migration, not to test a causal hypothesis; the experiments reported
below are sources of evidence inside the case.

\section{Goal and Research Questions}
\label{sec:goal-rqs}
The goal stated in Chapter 1 is refined into five research questions, one per metric group.
\begin{description}[leftmargin=1cm]
\item[RQ1 (Effort)] What steps, decisions and effort were required to derive the ontology from the data contracts?
\item[RQ2 (Fidelity)] To what extent are the constructs of the data contracts represented in the ontology as typed knowledge, as opaque literals, or not at all?
\item[RQ3 (Routing effectiveness)] How does the accuracy of the server's dataset router change when it is fed by the contracts, by the ontology, or by both?
\item[RQ4 (Agent effectiveness)] How does the accuracy of an LLM-based agent in choosing the dataset change under the same conditions?
\item[RQ5 (Cost)] How much context does the agent consume under each condition?
\end{description}

\section{Context and Case Selection}
The case is TIAGO Dados Abertos, an open-source MCP server maintained by the Brazilian
power system operator (ONS). At version 1.1.0 it describes 80 open datasets with 80 data
contracts written in the Open Data Contract Standard, with 887 columns and 351 governed
metrics. It exposes four tools (\texttt{buscar\_dataset}, \texttt{listar\_datasets},
\texttt{descrever\_dataset}, \texttt{executar\_sql}); dataset routing combines BM25 with
keyword rules over the contracts.

The case was selected for three reasons. The contracts are unusually rich: besides schema,
each one carries a \texttt{semantics} block with row grain, metric definitions and
aggregation rules. The domain has quantities that cannot be aggregated in the same way
(average power in MWmed, stored energy in MWmes), so a wrong answer can be silently
plausible. And the author works at the organisation and had access to question sets written
by its staff. \todo{Describe the author's role; add the architecture figure required by
item 5 of the manual.}

\section{Data Collection}
\textbf{Contracts and ontology.} The ontology was built twice. A first, exploratory attempt
with an AI coding assistant produced a small ontology with hand-curated aliases. It is kept
as a tagged snapshot and is not used in the results. The second attempt, reported here,
separates a hand-written schema (T-Box, 168 lines of Turtle) from generated content (A-Box).
The schema reuses DCAT and Dublin Core for dataset metadata, RDF Data Cube for measures and
dimensions, CSVW for columns, SKOS for controlled vocabularies and QUDT for quantity kinds
and units. A generator reads the 80 contracts and emits 39{,}390 triples; no instance data
is written by hand. Every construct occurrence found in a contract is recorded in a ledger
with its level of representation.

\textbf{Question sets.} Two sets written by ONS staff for an earlier question-answering
system were used: a curated test battery with 133 questions, of which 106 are usable for
routing (27 follow-up questions that depend on a previous turn and rows with no expected
dataset were excluded), and a selector benchmark with 204 questions that have exactly one
correct dataset. Neither set repeats a question indexed by the server. \todo{Record the
authorisation from ONS and the provenance of the 204 set.}

\textbf{Process records.} For the first attempt, the conversation log of the assistant and
the file snapshots it kept; for the second, the version history of the study repository.

\section{Analysis}
For RQ2, construct occurrences are classified as \emph{semantic} (a typed node or controlled
concept linked in the graph), \emph{literal} (content preserved as opaque text, SQL or JSON),
\emph{dropped}, or not applicable when the contract explicitly declares no value. The share
carried by standard vocabularies is also reported.

For RQ3, the production router is never modified; only the catalogue it receives changes.
Five catalogues are compared: schema only, the contracts, the contracts with an ontology
re-rank, and two catalogues rebuilt from the ontology graph, one with semantic content only
and one with semantic and literal content. The re-rank moves members of a measure family at
the grain named in the question ahead of their siblings and has no tunable weight. Top-1,
top-3 and MRR are reported with 95\% Wilson intervals.

For RQ4 and RQ5, a Claude Sonnet 5.5 model on Amazon Bedrock receives a search tool, a
describe tool and a final-answer tool, all backed by the catalogue of the condition. The
answer is correct when the chosen dataset is one of the valid ones. Input tokens per
question measure the context cost. Sampling parameters cannot be set on this model, so the
round is not deterministic; each condition was run once.

Paired differences use the exact McNemar test with Holm correction, and paired bootstrap
intervals with 10{,}000 resamples.

\section{Results}

\subsection{RQ1: Effort}
Table~\ref{tab:effort} summarises the two attempts. In the first, the assistant log holds
19 user prompts, 204 tool calls and 46 minutes of assistant time over four working days,
with 73 intermediate versions of 26 files. Its outcome covered three contract constructs,
and the phenomenon aliases were hand-tuned until coverage rose from 22 to 74 of 77 datasets.
In the second, the mapping is a program: a generator of about 400 lines of Python and 168
lines of hand-written schema; the whole study code has about 1{,}600 lines.
\todo{Add the person-time of the second attempt and the list of design decisions.}

\begin{table}[ht]
\centering\small
\caption{Effort of the two attempts at deriving the ontology.}
\label{tab:effort}
\begin{tabular}{lrr}
\toprule
 & First attempt & Second attempt \\
\midrule
Contract construct types represented & 3 of 65 & 59 of 65 \\
Hand-written instance data & partly & none \\
Instance data & 77 dataset bindings & 39{,}390 triples \\
Schema, hand-written lines & 254 & 168 \\
File versions kept & 73 snapshots & \todo{n} commits \\
\bottomrule
\end{tabular}
\end{table}

\subsection{RQ2: Fidelity}
Table~\ref{tab:fidelity} shows what happened to the 26{,}839 construct occurrences found in
the 80 contracts; 484 further occurrences are explicit ``no value'' markers. More than half
of the content became typed knowledge and 40\% was preserved as literals. What was dropped
belongs to the physical layer and to contract housekeeping.

\begin{table}[ht]
\centering\small
\caption{Construct occurrences of the data contracts by level of representation in the ontology.}
\label{tab:fidelity}
\begin{tabular}{lrrr}
\toprule
Contract block & Semantic & Literal & Dropped \\
\midrule
Catalogue and identity & 2{,}063 & 240 & 0 \\
Grain, dimensions and classification & 1{,}304 & 164 & 4 \\
Metrics & 4{,}045 & 2{,}357 & 0 \\
Columns & 5{,}758 & 5{,}513 & 0 \\
Value aliases & 1{,}606 & 0 & 16 \\
Agent support (questions and SQL) & 413 & 1{,}217 & 0 \\
Physical layer and operations & 0 & 1{,}297 & 842 \\
\midrule
Total & 15{,}189 (56.6\%) & 10{,}788 (40.2\%) & 862 (3.2\%) \\
\bottomrule
\end{tabular}
\end{table}

Of the 25{,}977 represented occurrences, 13{,}480 (51.9\%) are carried by a standard
vocabulary. Units show the limit of reuse: only 14 of the 45 units found in the contracts
have a QUDT equivalent, and the units specific to the Brazilian sector (MWmed, MWmes,
R\$/MWh) do not. The migration also surfaced inconsistencies that the contracts tolerate:
the same domain written with and without accents, two spellings of the same unit, and two
names for the same aggregation.

\subsection{RQ3: Routing effectiveness}
\begin{table}[ht]
\centering\small
\caption{Top-1 accuracy of the server's router (95\% Wilson interval).}
\label{tab:routing}
\begin{tabular}{lcc}
\toprule
Catalogue given to the router & Battery ($n=106$) & Selector set ($n=204$) \\
\midrule
Schema only & 0.330 [0.25, 0.42] & 0.382 [0.32, 0.45] \\
Contracts (production) & 0.679 [0.59, 0.76] & 0.750 [0.69, 0.80] \\
Contracts + ontology re-rank & 0.689 [0.60, 0.77] & 0.745 [0.68, 0.80] \\
Ontology, semantic content only & 0.632 [0.54, 0.72] & 0.667 [0.60, 0.73] \\
Ontology, semantic and literal content & 0.679 [0.59, 0.76] & 0.750 [0.69, 0.80] \\
\bottomrule
\end{tabular}
\end{table}

The catalogue rebuilt from the full ontology reproduces the contracts: no question is routed
differently in either set. With semantic content only, accuracy drops by 5 to 8 points; the
difference is significant on the larger set (Holm-adjusted $p=0.035$) and not on the battery
($p=0.44$). The re-rank changes one answer in each set, in opposite directions. On the
questions rated hard in the battery, accuracy is at most 0.40 under every condition.

\subsection{RQ4: Agent effectiveness}
\begin{table}[ht]
\centering\small
\caption{Accuracy of the LLM-based agent in choosing the dataset (95\% Wilson interval) and
mean input tokens per question.}
\label{tab:agent}
\begin{tabular}{lcrcr}
\toprule
 & \multicolumn{2}{c}{Battery ($n=106$)} & \multicolumn{2}{c}{Selector set ($n=204$)} \\
Catalogue behind the tools & Accuracy & Tokens & Accuracy & Tokens \\
\midrule
Schema only & 0.660 [0.57, 0.74] & 3{,}999 & 0.838 [0.78, 0.88] & 3{,}624 \\
Contracts & 0.858 [0.78, 0.91] & 6{,}777 & 0.912 [0.87, 0.94] & 6{,}117 \\
Ontology, semantic only & 0.868 [0.79, 0.92] & 4{,}460 & 0.917 [0.87, 0.95] & 4{,}100 \\
Ontology, semantic and literal & 0.877 [0.80, 0.93] & 5{,}016 & 0.912 [0.87, 0.94] & 4{,}848 \\
\bottomrule
\end{tabular}
\end{table}

The agent is far more accurate than the router alone, and the semantic layers still matter:
contracts add 20 points over schema on the battery and 7 on the selector set (Holm-adjusted
$p<0.01$ in both). Between contracts and either ontology catalogue there is no detectable
difference in accuracy: at most seven discordant questions in any comparison, and no adjusted $p$ below 1. With
schema only, the agent gives no answer for 12\% of the battery questions; with any semantic
layer, for at most 2\%.

\subsection{RQ5: Cost}
The semantic-only ontology catalogue delivers the accuracy of the contracts with about a
third less context: 2{,}317 fewer input tokens per question on the battery (95\% interval
1{,}145 to 3{,}587) and 2{,}017 fewer on the selector set (1{,}199 to 2{,}913). The whole
round consumed 6.0 million input tokens and 0.23 million output tokens.

\section{Discussion and Lessons Learned}
\textbf{A generated ontology is a lossless carrier, so its value lies elsewhere.} Because
every routing-relevant construct was migrated, the full ontology behaves exactly like the
contracts. An ontology derived from complete contracts cannot beat them on a task the
contracts already support; it can only match them, reorganise them, or expose what they lack.

\textbf{Vocabulary matters more than structure.} The first attempt kept three structural
axes and dropped the vocabulary, and its ontology-only condition routed worse than a plain
schema. The second attempt carries labels and aliases as typed annotations and routes close
to the contracts without any free text.

\textbf{An agent needs less text than a lexical router.} Free text and example questions are
worth 5 to 8 points to the BM25 router and nothing measurable to the agent, which reaches the
same accuracy on typed content with a third less context.

\textbf{The migration works as an audit of the contracts.} Normalising values exposed
inconsistent spellings, and typing the grain exposed contracts that declare no entity where
one exists (the equivalent-reservoir datasets are typed at subsystem grain).

\textbf{The remaining errors are confusions between siblings.} With any semantic layer the
agent misses 11 of the 106 battery questions and 15 of the 204 selector questions under
every condition. The wrong choices are almost always a neighbouring dataset of the same
subject: daily load instead of monthly load, the general instead of the detailed scheduling
balance, plant generation instead of installed capacity. Two selector questions expect a
dataset that the catalogue marks as discontinued and therefore never offers.

\textbf{Cross-dataset knowledge did not help routing.} The family-grain re-rank, the one
mechanism that uses knowledge absent from individual contracts, was neutral.

\section{Threats to Validity}
\textbf{Construct validity.} Choosing the right dataset is a necessary step, not the final
answer; neither the SQL nor the returned value was evaluated. The fidelity levels depend on
mapping decisions made by the author, which are explicit in the generator. The questions
were written for another system and their expected datasets reflect that system's catalogue.

\textbf{Internal validity.} The ontology content comes from the same contracts it is
compared with, which bounds what it can add. The agent round was run once per condition on a
non-deterministic model. The question sets had already been used to evaluate the first attempt when the second
ontology was designed: the generator has no parameter fitted to them, but the design was
not blind to those results.

\textbf{External validity.} One organisation, one catalogue of 80 datasets, questions in
Portuguese, one router and one model. The contracts are richer than usual, which favours the
contracts condition and may not hold elsewhere.

\textbf{Reliability.} Code, ontology, ledger and aggregate results are public in the study
repository, pinned to a corpus commit. The question sets and the assistant log are not
public; the scripts that convert them are. The names of the IEC CIM classes used in the
alignments were checked against the CIM17 schema documentation; the namespace IRI was not. Much of the code was produced with AI assistants under the author's
direction.

\section{Concluding Remarks}
\todo{Two or three sentences closing the chapter.}
